% Laporan: pipeline analisis sentimen (Gemini LLM + SVM) — Jupyter Notebook
% Kompilasi: pdflatex (2x) atau latexmk -pdf
% Ganti cuplikan kode: simpan screenshot per sel di folder figures/ atau gunakan \includegraphics

\documentclass[11pt,a4paper]{article}

\usepackage[a4paper,margin=2.5cm]{geometry}
\usepackage[indonesian]{babel}
\usepackage{graphicx}
\usepackage{hyperref}
\usepackage{enumitem}
\usepackage{caption}
\usepackage{booktabs}
\usepackage{amsmath}
\usepackage{xcolor}

\hypersetup{colorlinks=true, linkcolor=blue, urlcolor=blue}

% --- Cuplikan kode: kotak placeholder (hapus dan ganti \includegraphics jika sudah ada gambar) ---
\newcommand{\CuplikanKode}[2]{%
  \begin{figure}[htbp]
    \centering
    % \includegraphics[width=0.95\textwidth]{figures/#1}
    \fbox{%
      \parbox{0.92\textwidth}{%
        \centering\vspace{3.2cm}
        \textcolor{gray}{\textit{Tempat screenshot / cuplikan kode: #2}}\\
        \vspace{0.3cm}
        \textcolor{gray}{\small Berkas gambar disarankan: \texttt{figures/#1}}
        \vspace{3.2cm}%
      }%
    }
    \caption{#2}
    \label{fig:#1}
  \end{figure}%
}

\title{\textbf{Laporan Teknis:\\Analisis Sentimen Fitur SpayLater (ShopeePay)}\\[0.4em]
\large Autolabeling dengan Google Gemini \& Klasifikasi dengan SVM}
\author{%
  \textit{(Nama penulis / NIM)}\\
  Notebook sumber: \texttt{spaylater\_sentiment\_analysis\_gemini.ipynb}%
}
\date{\today}

\begin{document}
\maketitle

\begin{abstract}
\noindent
Laporan ini mendokumentasikan alur kerja pada notebook Jupyter untuk analisis sentimen teks media sosial terkait SpayLater/ShopeePay: prapemrosesan bahasa Indonesia (Sastrawi), pelabelan otomatis tiga kelas (\textit{positif}, \textit{netral}, \textit{negatif}) memakai API Gemini, pelatihan model \textit{Support Vector Machine} (SVM) dengan TF--IDF, evaluasi ketepatan pada data uji, serta visualisasi \textit{wordcloud}. Setiap bagian kode dijelaskan dengan asumsi satu sel notebook setara dengan satu cuplikan layar kode.
\end{abstract}

\tableofcontents
\newpage

%==============================================================================
\section{Penjelasan per sel kode (cuplikan notebook)}
\label{sec:sel-kode}

\noindent
\textbf{Catatan:} Nomor sel mengikuti indeks nol (\texttt{Cell 0}, \texttt{Cell 1}, \ldots) seperti urutan di \texttt{spaylater\_sentiment\_analysis\_gemini.ipynb}. Hanya sel bertipe \textbf{kode} yang diberi cuplikan; sel Markdown berisi narasi di notebook dan tidak diduplikasi di sini sebagai gambar.

\subsection*{Sel 3 --- Daftar dependensi (\texttt{requirements.txt})}
\CuplikanKode{sel03_requirements}{Sel 3: menulis berkas \texttt{requirements.txt} berisi \texttt{google-generativeai}, Sastrawi, pandas, scikit-learn, matplotlib, seaborn, dan wordcloud agar lingkungan dapat direplikasi.}

\subsection*{Sel 4 --- Instalasi paket dari \texttt{requirements.txt}}
\CuplikanKode{sel04_pip_requirements}{Sel 4: perintah \texttt{pip install -r requirements.txt} untuk menginstal dependensi utama analisis.}

\subsection*{Sel 5 --- Paket unduhan Mega}
\CuplikanKode{sel05_pip_mega}{Sel 5: menginstal \texttt{mega.py} untuk mengunduh dataset dari tautan Mega.nz (opsional jika data sudah lokal).}

\subsection*{Sel 6 --- Penyesuaian \texttt{tenacity}}
\CuplikanKode{sel06_pip_tenacity}{Sel 6: memutakhirkan \texttt{tenacity} guna mengurangi konflik dependensi dengan \texttt{mega.py} atau paket lain di lingkungan yang sama.}

\subsection*{Sel 7 --- Unduh dataset dari Mega}
\CuplikanKode{sel07_mega_download}{Sel 7: login anonim Mega, mengunduh berkas Excel/CSV proyek dari URL yang ditentukan, dan mencetak path berkas hasil unduhan.}

\subsection*{Sel 8 --- Kunci API Gemini}
\CuplikanKode{sel08_api_key}{Sel 8: menetapkan variabel lingkungan \texttt{GEMINI\_API\_KEY} untuk autentikasi ke Google AI Studio / Gemini API. \textbf{Jangan} menyertakan kunci asli dalam laporan versi akhir; gunakan variabel lingkungan atau redaksi.}

\subsection*{Sel 10 --- Impor pustaka dan konfigurasi global}
\CuplikanKode{sel10_config}{Sel 10: impor NumPy, pandas, scikit-learn, matplotlib, seaborn, WordCloud; konfigurasi path data (\texttt{CSV\_PATH}), folder ekspor (\texttt{EXPORT\_DIR}, \texttt{RAW\_RESPONSE\_DIR}), batas batch LLM (\texttt{MAX\_TOKEN\_DATA}, \texttt{MAX\_LLM\_BATCH\_ITEMS}), model Gemini (\texttt{LLM\_MODEL}), token keluaran, dan jeda antar batch (\texttt{LLM\_BATCH\_SLEEP\_SEC}).}

\subsection*{Sel 11 --- Estimasi token (\texttt{count\_tokens})}
\CuplikanKode{sel11_tiktoken}{Sel 11: mendefinisikan \texttt{count\_tokens} memakai \texttt{tiktoken} (\texttt{cl100k\_base}) jika tersedia; jika tidak, perkiraan kasar $\max(1, \lfloor |s|/4 \rfloor)$ untuk membantu membagi batch LLM.}

\subsection*{Sel 13 --- Muat dataset}
\CuplikanKode{sel13_read_excel}{Sel 13: membaca Excel dengan pandas, menampilkan cuplikan baris awal, dan memastikan kolom \texttt{full\_text} ada.}

\subsection*{Sel 14 --- Sampel acak (opsional)}
\CuplikanKode{sel14_sample}{Sel 14: membuat sampel acak berukuran \texttt{LEN\_SAMPLE\_DATA} untuk uji token; baris \texttt{df = df\_sampled} dapat diaktifkan untuk membatasi ukuran data.}

\subsection*{Sel 15 --- Total token pada seluruh baris}
\CuplikanKode{sel15_token_sum}{Sel 15: menghitung \texttt{token\_count} per baris pada \texttt{full\_text} dan menjumlahkan total token korpus.}

\subsection*{Sel 16 --- Rasio estimasi token}
\CuplikanKode{sel16_ratio}{Sel 16: operasi aritmetika sederhana (mis. perbandingan total token terhadap sampel) untuk estimasi skala data.}

\subsection*{Sel 17 --- Jumlah baris}
\CuplikanKode{sel17_len}{Sel 17: \texttt{len(df)} untuk mencetak jumlah observasi.}

\subsection*{Sel 18 --- Histogram panjang token}
\CuplikanKode{sel18_hist}{Sel 18: histogram distribusi \texttt{token\_count} untuk melihat sebaran panjang teks.}

\subsection*{Sel 20 --- Fungsi prapemrosesan teks}
\CuplikanKode{sel20_preprocess}{Sel 20: kamus singkatan (\texttt{SLANG\_MAP}), regex URL/email/mention/hashtag, normalisasi Unicode, stemming dan stopword Sastrawi, serta \texttt{preprocess\_series} untuk diterapkan ke kolom DataFrame.}

\subsection*{Sel 22 --- Demo langkah prapemrosesan}
\CuplikanKode{sel22_demo}{Sel 22: fungsi \texttt{demo\_steps} mencetak tiap tahap prapemrosesan untuk beberapa sampel pertama (\texttt{full\_text}) guna verifikasi visual.}

\subsection*{Sel 24 --- Terapkan prapemrosesan ke seluruh data}
\CuplikanKode{sel24_text_clean}{Sel 24: membuat kolom \texttt{text\_clean} dari \texttt{full\_text}; kolom ini dipakai sebagai fitur masukan SVM.}

\subsection*{Sel 25 --- Penanda / sel cadangan}
\CuplikanKode{sel25_placeholder}{Sel 25: sel kode kosong atau komentar penanda; dapat diabaikan atau diisi eksperimen singkat.}

\subsection*{Sel 27 --- Autolabel dengan Gemini (inti)}
\CuplikanKode{sel27_autolabel}{Sel 27: integrasi \texttt{google.generativeai}; normalisasi label; parsing JSON array respons LLM (\texttt{parse\_llm\_sentiment\_labels\_json}, \texttt{extract\_first\_json\_array}); pembagian batch; penyimpanan respons mentah ke \texttt{raw\_response/}; rekursi pemecahan batch jika parsing gagal; penanganan kuota; pengisian kolom \texttt{label} pada DataFrame.}

\subsection*{Sel 29 --- Nilai unik label}
\CuplikanKode{sel29_label_unique}{Sel 29: \texttt{df['label'].unique()} untuk memeriksa kelas setelah autolabel (atau jika label sudah ada di file).}

\subsection*{Sel 30 --- Pelatihan SVM dan pencarian hyperparameter}
\CuplikanKode{sel30_svm_grid}{Sel 30: pemetaan label ke integer; \texttt{Pipeline} TF--IDF (1--2 gram) + \texttt{SVC} linear berbobot kelas seimbang; \texttt{GridSearchCV} pada kombinasi \texttt{max\_features} dan \texttt{C}; pencarian \texttt{test\_size} terbaik pada hold-out terstratifikasi; menyimpan model dan data uji terbaik.}

\subsection*{Sel 32 --- Matriks kebingungan dan laporan klasifikasi}
\CuplikanKode{sel32_confusion}{Sel 32: prediksi pada \texttt{X\_test}; matriks kebingungan dengan heatmap Seaborn; cetak \texttt{classification\_report} (precision, recall, F1); simpan gambar dan teks laporan ke \texttt{EXPORT\_DIR}.}

\subsection*{Sel 34 --- Wordcloud}
\CuplikanKode{sel34_wordcloud}{Sel 34: menggabungkan \texttt{text\_clean} menjadi satu string, membuat \textit{wordcloud}, menampilkan dan menyimpan PNG.}

\subsection*{Sel 36 --- Ekspor artefak dan ZIP}
\CuplikanKode{sel36_export_zip}{Sel 36: menyimpan CSV berlabel, sampel prapemrosesan, \texttt{summary.json} (akurasi hold-out, parameter terbaik, ukuran data), mengompres isi folder ekspor menjadi satu berkas ZIP.}

\subsection*{Sel 37 --- Unduh ZIP (Google Colab)}
\CuplikanKode{sel37_colab_download}{Sel 37: \texttt{files.download} Colab untuk mengunduh arsip hasil; di lingkungan lokal, langkah ini dapat diganti membuka folder ekspor secara manual.}

\newpage
%==============================================================================
\section{Hasil analisis sentimen}
\label{sec:hasil}

\subsection{Ringkasan alur dan interpretasi umum}

Setelah seluruh pipeline dijalankan, \textbf{hasil analisis sentimen} dapat dibahas dalam kerangka berikut.

\begin{enumerate}[leftmargin=*]
  \item \textbf{Distribusi label (autolabel LLM).}
  Kolom \texttt{label} merepresentasikan sentimen tiap dokumen: \textit{positif}, \textit{netral}, atau \textit{negatif}.
  Sebaiknya dilaporkan jumlah dan proporsi per kelas (bisa dihitung dengan \texttt{df['label'].value\_counts(normalize=True)}) untuk melihat apakah data \textit{imbalance} dan apakah autolabel menghasilkan pola yang masuk akal untuk domain SpayLater/ShopeePay.

  \item \textbf{Pemilihan pembagian data dan model.}
  Notebook membandingkan beberapa nilai \texttt{test\_size} dan memilih kombinasi yang memberi \textbf{akurasi hold-out tertinggi} setelah \texttt{GridSearchCV} pada data latih.
  Pada \texttt{summary.json} (dan keluaran Sel 30) tercatat \texttt{best\_test\_size}, \texttt{holdout\_accuracy}, \texttt{best\_cv\_accuracy}, dan \texttt{best\_params} --- nilai inilah yang harus Anda kutip sebagai \textbf{hasil utama} kinerja SVM pada skenario data Anda.

  \item \textbf{Matriks kebingungan.}
  Heatmap membandingkan label aktual (baris) dan prediksi (kolom).
  Elemen diagonal besar berarti klasifikasi benar untuk kelas tersebut; kesalahan dominan di luar diagonal mengindikasikan pasangan kelas yang sering tertukar (misalnya netral vs positif), yang relevan untuk membahas kualitas teks atau ambiguitas opini pengguna.

  \item \textbf{Laporan klasifikasi (precision, recall, F1).}
  Untuk tiap kelas, \textit{precision} menjawab seberapa dapat dipercaya prediksi positif untuk kelas itu; \textit{recall} seberapa banyak kasus aktual yang tertangkap; \textit{F1} menyeimbangkan keduanya.
  Pada dataset tidak seimbang, akurasi global saja bisa menyesatkan; F1 per kelas dan rata-rata makro/terbobot lebih informatif.

  \item \textbf{Wordcloud.}
  Visual ini menonjolkan kata-kata frekuensi tinggi pada \texttt{text\_clean} setelah normalisasi.
  Kata yang dominan dapat mendukung narasi apakah wacana publik lebih sering menyentuh tema bunga, limit, keterlambatan, promo, dan sebagainya --- diselaraskan dengan distribusi sentimen.

  \item \textbf{Keterbatasan.}
  Label dari LLM bergantung pada kualitas prompt dan kuota API; SVM pada TF--IDF memodelkan asosiasi leksikal tanpa konteks mendalam seperti model besar.
  Hasil bersifat \textit{indikatif} untuk penelitian/exploratory analysis dan sebaiknya divalidasi dengan sampel anotasi manusia bila digunakan untuk keputusan penting.
\end{enumerate}

\subsection{Tabel ringkasan (isi dari notebook Anda)}

\noindent
Isi tabel berikut dengan angka yang muncul di keluaran Sel 30--32 / berkas \texttt{summary.json} setelah Anda menjalankan notebook.

\begin{table}[htbp]
  \centering
  \caption{Ringkasan metrik evaluasi (lengkapi setelah eksekusi notebook)}
  \begin{tabular}{@{}ll@{}}
    \toprule
    \textbf{Ukuran} & \textbf{Nilai} \\
    \midrule
    Jumlah sampel (\texttt{n\_samples}) & \ldots \\
    Proporsi kelas negatif / netral / positif & \ldots \\
    \texttt{test\_size} terpilih & \ldots \\
    Akurasi hold-out & \ldots \\
    Akurasi CV terbaik (latih) & \ldots \\
    Parameter terbaik (\texttt{max\_features}, \texttt{C}) & \ldots \\
    \bottomrule
  \end{tabular}
\end{table}

\subsection{Paragraf contoh pembahasan (sesuaikan angka)}

\textit{Model SVM dengan vektorisasi TF--IDF pada teks yang telah dinormalisasi mencapai akurasi hold-out sebesar \textbf{[\ldots]} dengan pembagian uji \textbf{[\ldots\%]}. Dari matriks kebingungan, kelas yang paling sulit dipisahkan adalah \textbf{[\ldots]}. Wordcloud menunjukkan kata-kata seperti \textbf{[\ldots]} mendominasi korpus, selaras dengan proporsi sentimen \textbf{[\ldots]}. Autolabel Gemini mempercepat pelabelan skala besar; namun sampel validasi manual disarankan pada kasus ambigu.}

\vspace{1em}
\noindent
\textbf{Referensi kode:} seluruh implementasi terperinci berada di \texttt{spaylater\_sentiment\_analysis\_gemini.ipynb}.

\end{document}
